% Options for packages loaded elsewhere
\PassOptionsToPackage{unicode}{hyperref}
\PassOptionsToPackage{hyphens}{url}
\PassOptionsToPackage{dvipsnames,svgnames,x11names}{xcolor}
%
\documentclass[
  10pt,
]{article}
\usepackage{amsmath,amssymb}
\usepackage{iftex}
\ifPDFTeX
  \usepackage[T1]{fontenc}
  \usepackage[utf8]{inputenc}
  \usepackage{textcomp} % provide euro and other symbols
\else % if luatex or xetex
  \usepackage{unicode-math} % this also loads fontspec
  \defaultfontfeatures{Scale=MatchLowercase}
  \defaultfontfeatures[\rmfamily]{Ligatures=TeX,Scale=1}
\fi
\usepackage{lmodern}
\ifPDFTeX\else
  % xetex/luatex font selection
\fi
% Use upquote if available, for straight quotes in verbatim environments
\IfFileExists{upquote.sty}{\usepackage{upquote}}{}
\IfFileExists{microtype.sty}{% use microtype if available
  \usepackage[]{microtype}
  \UseMicrotypeSet[protrusion]{basicmath} % disable protrusion for tt fonts
}{}
\makeatletter
\@ifundefined{KOMAClassName}{% if non-KOMA class
  \IfFileExists{parskip.sty}{%
    \usepackage{parskip}
  }{% else
    \setlength{\parindent}{0pt}
    \setlength{\parskip}{6pt plus 2pt minus 1pt}}
}{% if KOMA class
  \KOMAoptions{parskip=half}}
\makeatother
\usepackage{xcolor}
\usepackage[margin=25mm]{geometry}
\usepackage{longtable,booktabs,array}
\usepackage{calc} % for calculating minipage widths
% Correct order of tables after \paragraph or \subparagraph
\usepackage{etoolbox}
\makeatletter
\patchcmd\longtable{\par}{\if@noskipsec\mbox{}\fi\par}{}{}
\makeatother
% Allow footnotes in longtable head/foot
\IfFileExists{footnotehyper.sty}{\usepackage{footnotehyper}}{\usepackage{footnote}}
\makesavenoteenv{longtable}
\setlength{\emergencystretch}{3em} % prevent overfull lines
\providecommand{\tightlist}{%
  \setlength{\itemsep}{0pt}\setlength{\parskip}{0pt}}
\setcounter{secnumdepth}{5}
\ifLuaTeX
\usepackage[bidi=basic]{babel}
\else
\usepackage[bidi=default]{babel}
\fi
\babelprovide[main,import]{american}
% get rid of language-specific shorthands (see #6817):
\let\LanguageShortHands\languageshorthands
\def\languageshorthands#1{}
\ifLuaTeX
  \usepackage{selnolig}  % disable illegal ligatures
\fi
\usepackage{bookmark}
\IfFileExists{xurl.sty}{\usepackage{xurl}}{} % add URL line breaks if available
\urlstyle{same}
\hypersetup{
  pdftitle={Temporal Attention Neurons: Separating History Sensitivity, Addressing, and Recoverable Memory},
  pdfauthor={Li Zexu},
  pdflang={en-US},
  colorlinks=true,
  linkcolor={Maroon},
  filecolor={Maroon},
  citecolor={Blue},
  urlcolor={Blue},
  pdfcreator={LaTeX via pandoc}}

\title{Temporal Attention Neurons: Separating History Sensitivity,
Addressing, and Recoverable Memory}
\author{Li Zexu}
\date{3 November 2026 --- research manuscript draft}

\begin{document}
\maketitle

\section{Abstract}\label{abstract}

A temporal attention mechanism can change a neuron's response to recent
history without making that history recoverable through its permitted
output. We synthesize the recorded Temporal Attention Neuron (TAN)
research into an evidence-based account of this distinction. The
evidence spans static addressing, native finite-window dynamics,
task-specific response collisions, a controlled readout-training
intervention, and comparisons with ordinary adaptive filters. Static
scalar metric addressing correctly identifies all 6,705 intermediate-key
queries in the archived noiseless three-key condition, whereas the
signed scalar-product rule identifies none. This positive result
presupposes an available key--value array. In a separate native temporal
task, 545 target--distractor swap pairs with different labels have
identical complete allowed responses, despite different ordered internal
windows. A later clean-only training intervention improves mean clean
balanced accuracy from 32.5881\% to 38.1199\%, but does not eliminate
this structural limitation. In a matched-calibration detection
experiment, TAN's balanced accuracy exceeds an ordinary adaptive
comparator under independent noise but falls below it under stationary
AR(0.9) noise. We distinguish these locally executed records from
exploratory Chat-sandbox constructions that have not been independently
replicated here. The resulting contribution is a task- and
resource-specific account of what the architecture computes, where
information becomes inaccessible, and which apparent gains depend on
readout, training composition, or additional mechanisms. No universal
superiority, biological implementation, or independent experimental
replication is established.

\section{Introduction}\label{introduction}

Attention is a computation over a representation, not by itself a
guarantee that the representation retains the information required by a
task. The Transformer makes content-dependent weighting operational at
network scale {[}1{]}. TAN asks a smaller architectural question: what
can a discrete dynamical unit do when its membrane update receives a
surprise-gated, attention-weighted summary of recent inputs? Answering
this question requires separating at least four properties: history
sensitivity, address selection, preservation of a key--value
association, and recovery through a specified observation.

The original project developed the TAN model, its experimental tasks,
and an architectural progression from state and geometry to routing,
binding, and composition. Existing manuscripts already discuss these
distinctions. This paper is a retrospective synthesis of the subsequent
evidence, not a duplicate presentation of those manuscripts or a new
experimental campaign. In particular, later evidence qualifies a simple
monotonic capability ladder: a representation may contain a target while
its output cannot expose it; a scalar metric can address an existing
array without solving dynamic storage; and a training change can improve
observed accuracy without repairing an information-losing map.

We organize the evidence around three questions. Which positive results
survive an explicit account of inputs, state, and readout? Which
negative results are mathematical obstructions within a fixed contract,
rather than broad claims about neurons? Which performance differences
remain when ordinary baselines and training composition are considered?
The answer is a set of bounded capabilities, not a single success or
failure label for all implementations called TAN.

\section{Materials and evidence
provenance}\label{materials-and-evidence-provenance}

\subsection{Source hierarchy and prior
manuscripts}\label{source-hierarchy-and-prior-manuscripts}

The original source baseline is commit
\texttt{85213f0d417d54e4f48b301ed1c0e4200ef3c421}. The screened research
publication is commit \texttt{222caa0201feff1e37a06049aa44c0e59e1cbaba},
containing 680 archived scientific files in four collections: audit v11
(440), native eight-stage studies (110), exploratory studies (100), and
the local clean-training control (30). The exact screened ZIP has
189,164,458 bytes and SHA-256
\texttt{ca8a2ae51b4e8f18ab341fa483ca9970a7a3077c5d4e4013e5df13a9811b6167}.
Hash agreement establishes file identity, not scientific validity.

An inventory made before drafting records 174 existing manuscript/report
sources. The principal earlier manuscript families are
\texttt{snapshot/manuscript}, \texttt{snapshot/paper}, and
\texttt{snapshot/paper2}. Later correction fragments under the
stage-localization audit are amendments, not an additional independent
study. We retain these sources unchanged and use the latest completed
report relevant to each claim. For static addressing, the authoritative
experiment is P1A Rev2, not the superseded first run. For clean
training, the frozen completed protocol supersedes an earlier proposal
whose prospective wording remains historical.

During final publication checks, main advanced to commit
\texttt{65b0d2e9310a5790e653b77114bf4b6c7aa93e9c}, adding a TAN-III
manuscript focused on decoder floors, carrier rank, and a proposed
dendritic interpretation. We inspected that concurrent manuscript and
retain it unchanged. The present synthesis is a separate evidence-scoped
account centered on observable recovery and the training-composition
control. Shared numerical findings are not counted as additional
experiments. The proposed biological interpretation is not treated here
as an experimentally demonstrated implementation.

The evidence ledger accompanying this manuscript gives
repository-relative links and hashes for the sources actually used.
Reported numerical findings below are historical results checked against
retained evidence. Writing and building this paper does not constitute a
new run of the underlying experiments.

\subsection{Execution and attribution}\label{execution-and-attribution}

The model and original project are attributed to Li Zexu, consistent
with the existing repository. Subsequent audit design, coding,
mathematical analysis, computational checking, and this synthesis
involved AI assistance. This statement does not assign every later
derivation or implementation exclusively to the human author. AI systems
are not listed as authors; human review remains necessary before journal
submission.

We distinguish three evidence levels. First, exact arguments establish
consequences of explicitly stated mathematical maps. Second, the native
and local audit collections record executed computations, saved inputs,
outputs, and source provenance. Some include separately implemented
checkers or isolated replays; these are computational verification
within an AI-assisted workflow, not independent human replication.
Third, the \texttt{exploratory/chat-*} collections record Chat-sandbox
conceptual experiments, including same-agent checks and reruns. We have
inspected their documentation but have not independently reproduced
those experiments in the current environment. The paper is an unreviewed
research draft, not peer-reviewed evidence.

\section{Models, contracts, and
analysis}\label{models-contracts-and-analysis}

\subsection{Native finite-window TAN}\label{native-finite-window-tan}

For the W5 native model used in the temporal task, let \(H_t\) be the
ordered window containing the current input and four preceding inputs,
and let \(h_{t-1}\) be the previous post-reset membrane state. The
principal operations are

\[
\mu_t=\frac{1}{5}\sum_{z\in H_t}z,\qquad
S_t=\max(0,x_t-\mu_t),\qquad q_t=2S_t,
\]

\[
C_t=\sum_i\alpha_{ti}H_{ti},\qquad
A_t=\tanh(S_t)C_t,\qquad
u_t=0.5h_{t-1}+A_t,
\]

\[
y_t=\mathbf{1}\{u_t>0.5\},\qquad
h_t=(1-y_t)u_t.
\]

The attention weights use exponentials of \(q_tH_{ti}\), with a small
additive denominator constant. Implementation details matter. The native
v8 collision proof uses the historical unshifted exponential
normalization with \(\delta=10^{-9}\). The clean-training control uses
the source's stabilized softmax: subtract the maximum score before
exponentiation and then add the same denominator constant. Because of
that constant, these implementations are not exactly algebraically
identical. We do not silently transfer an exact floating-point result
between them {[}E2, E3{]}.

The full finite-window system is Markovian in \((h_t,H_t)\). Earlier
natural-collision results show that membrane alone is insufficient; they
do not establish non-Markovianity of the full state. The five window
values plus membrane are a 48-byte logical payload under binary64
storage, excluding object overhead, temporaries, and the external
decoder.

\subsection{Higher-moment sensitivity is not positional
addressing}\label{higher-moment-sensitivity-is-not-positional-addressing}

For the ideal normalization without \(\delta\), define

\[
Z(q)=\sum_i e^{qH_i},\qquad C(q)=\frac{d}{dq}\log Z(q).
\]

A local expansion at \(q=0\) gives

\[
C(q)=\mu+q\kappa_2+\frac{q^2}{2}\kappa_3+\cdots,
\]

where \(\kappa_j\) are cumulants of the empirical window distribution.
Thus native attention generally responds to information beyond mean and
variance. A clipped second-order approximation such as
\(\mu+q\operatorname{Var}(H)\) is an ablation, not an exact substitute.
This is an analytical property of exponential weighting, not evidence of
a novel statistical identity or of improved task accuracy.

The exact deployed normalization retains \(\delta\), so the ideal
expansion must not be quoted as an exact formula for its stored outputs.
Moreover, the context remains invariant to a permutation of a fixed
window when the query is unchanged. Sensitivity to the amplitude
distribution can coexist with loss of target position. The ordered
window can retain position while a pooled output discards it. This is
consistent with the broader distinction between set-invariant
computation and ordered representations {[}2{]}. The bundled upstream
control code contains a second-order context ablation, but its presence
alone does not establish a validated higher-moment performance advantage
{[}E3{]}.

\subsection{Static address selection}\label{static-address-selection}

P1A Rev2 grants the models the full stored key, value, and validity
arrays. For distinct keys and an exact query \(q=k_t\), the metric score

\[s_i=-(q-k_i)^2\]

has a unique maximum at the requested key. By contrast, a nonzero scalar
multiplier applied to ordered keys selects an extreme key; zero produces
a tie rather than unique intermediate-key identification. These are
elementary score properties. They do not prove a lower bound on neural
state dimension.

The vector-QK implementation's cosine score has the same hard ordering
as the metric on the observed domain, where \(|0.4(q-k)|<\pi\). Its
agreement with metric routing therefore has a direct explanation.
Soft-output comparisons require accounting for score curvature and
temperature: equal numeric temperature is not equal local
discrimination. Values and the final addition/subtraction operation are
explicitly available resources {[}E1{]}.

\subsection{Observation collision and training
control}\label{observation-collision-and-training-control}

In native Probe1, the target occurs at step 5, distractors at steps
6--8, query at step 9, and the continuation is zero. The only decoder
observation is

\[F=(u_9,u_{10},u_{11},u_{12},u_{13},y_9,y_{10},y_{11},y_{12},y_{13}).\]

The v8 study compares histories \((a,b,c,d,q)\) and \((b,a,c,d,q)\) with
\(a,c\in[0.615,0.635]\) and \(b,d,q\in[1.365,1.385]\). Their target
classes differ outside the original overlapping low/high amplitude
region. One anchor, all 32 box corners, and 512 frozen random draws give
545 paired cases {[}E2{]}.

The later training control asks a different question: whether the
composition of readout training data contributes to weak clean-task
accuracy. It holds the stabilized W5 source dynamics and 33-parameter
linear readout fixed. For each of three historical seeds, clean-only and
mixed arms each use 1,500 examples, balanced at 500 per class. The mixed
arm contains 250 clean and 250 distractor examples per class. Its 750
clean examples are a subset of the clean-only arm. Selection uses a
fixed permutation without consulting test values. Both arms evaluate the
same 1,001 clean examples per seed. Each arm uses its own train-only
standardization, zero initialization, L2 coefficient \(10^{-4}\), and
the frozen L-BFGS-B procedure, without validation-based tuning {[}E3{]}.

Balanced accuracy is the arithmetic mean of class recalls. Training
uncertainty is assessed only conditionally: 2,000 paired bootstrap draws
resample original test reuse-groups within each fixed seed, preserve
pairing between arms, and average equally across seeds. This does not
integrate over training-set, optimizer, or seed-selection uncertainty.

\subsection{Matched-calibration adaptive
baseline}\label{matched-calibration-adaptive-baseline}

The separate v11 detection task uses W10 TAN with signed surprise, a
magnitude gate, leak 0.8, strict threshold 1, and zero initial state. It
must not be confused with the W5 rectified-surprise classifier. The
ordinary comparator has two persistent scalars:

\[
m_t=0.8m_{t-1}+0.2x_t,\quad
I_t=\kappa\left|\frac98(x_t-m_t)\right|,\quad
u_t=0.8u_{t-1}^{\mathrm{post}}+I_t.
\]

Both models use the same threshold/reset and observation window. A
separate 512-history independent-noise calibration fixes
\(\kappa=0.840938900477691\) by matching pooled absolute effective
current, not accuracy or test outcomes. The comparator is an engineering
implementation of input adaptation, not a claimed reproduction of a
particular biological neuron. Adaptation-based separation of fast
signals and slower backgrounds already has primary-literature precedent
{[}3{]}.

The test contains 1,024 innovation histories paired across independent
and stationary AR(0.9) noise, step/catch conditions, and both models.
Each waveform lasts 60 updates, with step onset at 15. A detection is
any spike during updates 15--19; the readout therefore knows onset.
Misses receive capped delay 5. The bootstrap resamples whole innovation
groups, preserving every pairing. Calibration is frozen. Matching one
pooled current statistic does not match all state, computation, current
variance, or physical energy {[}E4{]}.

\section{Results}\label{results}

\subsection{Static addressing succeeds within its granted memory
interface}\label{static-addressing-succeeds-within-its-granted-memory-interface}

In the archived noiseless three-key condition, scalar metric attention
and vector QK each identify all 6,705 intermediate-key role queries,
whereas the signed scalar rule identifies zero. In the two noisy
conditions, both-address success is 9,876/10,000 and 7,986/10,000 for
each addressing method. These are observed task outcomes, not
unconditional robustness guarantees.

The retained verification recomputed 60,000 condition/trial records and
420,000 model-output rows from separately written score formulas. All
hard choices agreed exactly; the largest soft arithmetic discrepancy was
\(2.58\times10^{-14}\). The conditions do not supply 60,000 independent
histories because two noisy conditions share base histories and scaled
noise draws. This is saved-ledger verification, not regeneration of the
historical random experiment {[}E1{]}. The result supports static metric
addressing and challenges an overly narrow scalar-product obstruction.
It does not establish dynamic writing or compact key--value storage.

\subsection{Native response collisions separate stored from recoverable
content}\label{native-response-collisions-separate-stored-from-recoverable-content}

A rational interval certificate gives a step-8 gated current above
0.528613674287 throughout the specified box. Both native membranes
therefore reset. At query, the two windows contain the same multiset,
the same current input, and equal pooled context. The subsequent zero
continuation closes the gate, yielding the same response

\[F=(u_9,u_9/2,u_9/4,u_9/8,u_9/16,0,0,0,0,0).\]

At the anchor, \(u_9\approx0.33538361861604077\). All 545 native TAN
response pairs were also bit-identical in the recorded floating
implementation. The ordered windows nevertheless differed when queried
{[}E2{]}.

\begin{longtable}[]{@{}
  >{\raggedright\arraybackslash}p{(\columnwidth - 4\tabcolsep) * \real{0.3333}}
  >{\raggedleft\arraybackslash}p{(\columnwidth - 4\tabcolsep) * \real{0.3333}}
  >{\raggedleft\arraybackslash}p{(\columnwidth - 4\tabcolsep) * \real{0.3333}}@{}}
\toprule\noalign{}
\begin{minipage}[b]{\linewidth}\raggedright
Model or intervention
\end{minipage} & \begin{minipage}[b]{\linewidth}\raggedleft
Identical complete responses / 545 pairs
\end{minipage} & \begin{minipage}[b]{\linewidth}\raggedleft
Different spike responses / 545 pairs
\end{minipage} \\
\midrule\noalign{}
\endhead
\bottomrule\noalign{}
\endlastfoot
Native LIF & 545 & 0 \\
Native recency delay line & 0 & 0 \\
Native uniform gated context & 0 & 545 \\
Native TAN & 545 & 0 \\
TAN with step-8 voltage reset suppressed & 0 & 0 \\
\end{longtable}

Suppressing only the step-8 voltage reset restores analog separation but
not spike-only separation. It changes the dynamical rule and is a causal
control, not another successful run of unchanged TAN. The delay-line and
uniform-context results are valuable negative controls against a
universal advantage claim; they are not isolated single-component
interventions on TAN.

For an equal mixture of each opposite-label pair, any classifier of only
the common response incurs error at least one half. This is a restricted
indistinguishability statement. The corresponding conditional bound
under the ideal continuous original present-distractor generator is only
\(3.2\times10^{-9}\). The construction therefore cannot explain the
magnitude of average historical classification error. It proves a
failure mode, not its prevalence.

\subsection{Clean training improves the fixed readout, but only
modestly}\label{clean-training-improves-the-fixed-readout-but-only-modestly}

All six fitted optimizers converged. The recorded features exactly
matched the archived stabilized-source response. Table 2 reports
balanced accuracy on the shared clean test subsets {[}E3{]}.

\begin{longtable}[]{@{}
  >{\raggedright\arraybackslash}p{(\columnwidth - 6\tabcolsep) * \real{0.2500}}
  >{\raggedleft\arraybackslash}p{(\columnwidth - 6\tabcolsep) * \real{0.2500}}
  >{\raggedleft\arraybackslash}p{(\columnwidth - 6\tabcolsep) * \real{0.2500}}
  >{\raggedleft\arraybackslash}p{(\columnwidth - 6\tabcolsep) * \real{0.2500}}@{}}
\toprule\noalign{}
\begin{minipage}[b]{\linewidth}\raggedright
Historical seed
\end{minipage} & \begin{minipage}[b]{\linewidth}\raggedleft
Clean-only training (\%)
\end{minipage} & \begin{minipage}[b]{\linewidth}\raggedleft
Mixed training (\%)
\end{minipage} & \begin{minipage}[b]{\linewidth}\raggedleft
Difference (percentage points)
\end{minipage} \\
\midrule\noalign{}
\endhead
\bottomrule\noalign{}
\endlastfoot
20260904 & 38.5541 & 32.9776 & 5.5765 \\
20260905 & 38.0594 & 32.0142 & 6.0451 \\
20260906 & 37.7462 & 32.7725 & 4.9737 \\
Equal-seed mean & 38.1199 & 32.5881 & 5.5318 \\
\end{longtable}

The conditional 95\% percentile interval for the mean difference is
{[}3.3391, 7.6080{]} percentage points. All classes are predicted in
both arms, but changes are not uniform: for seed 20260905, low-class
recall under clean-only training is 11.51\%, while other recalls
improve. An aggregate improvement must not be described as improvement
for every class.

This intervention identifies training composition, including its induced
normalization, as a contributor under the frozen pipeline. It does not
isolate gradient conflict from preprocessing, identify a Bayes-optimal
decoder, or establish architectural superiority. Mean performance
remains modest, near the three-class balanced chance reference of one
third. Historical reused tests are not a new blind benchmark.

The clean-map question must also remain separate from the proven
distractor-swap collision. The retained v8 certificate contains nonzero
distractors. We did not locate a separate clean-only collision
certificate in the 680-file publication and do not relabel that theorem
as one. A poor clean linear classifier does not prove noninjectivity of
the clean input-to-response map; conversely, better clean training does
not prove that map is injective. Any additional clean-map collision
claim requires its exact source, witness, normalization, allowed query
and observation contract before it can be included as a verified result.

\subsection{Correlated noise reverses a detection
advantage}\label{correlated-noise-reverses-a-detection-advantage}

The complete primary detection outcomes are {[}E4{]}:

\begin{longtable}[]{@{}
  >{\raggedright\arraybackslash}p{(\columnwidth - 10\tabcolsep) * \real{0.1667}}
  >{\raggedright\arraybackslash}p{(\columnwidth - 10\tabcolsep) * \real{0.1667}}
  >{\raggedleft\arraybackslash}p{(\columnwidth - 10\tabcolsep) * \real{0.1667}}
  >{\raggedleft\arraybackslash}p{(\columnwidth - 10\tabcolsep) * \real{0.1667}}
  >{\raggedleft\arraybackslash}p{(\columnwidth - 10\tabcolsep) * \real{0.1667}}
  >{\raggedleft\arraybackslash}p{(\columnwidth - 10\tabcolsep) * \real{0.1667}}@{}}
\toprule\noalign{}
\begin{minipage}[b]{\linewidth}\raggedright
Noise
\end{minipage} & \begin{minipage}[b]{\linewidth}\raggedright
Model
\end{minipage} & \begin{minipage}[b]{\linewidth}\raggedleft
Misses / 1,024
\end{minipage} & \begin{minipage}[b]{\linewidth}\raggedleft
Catch alarms / 1,024
\end{minipage} & \begin{minipage}[b]{\linewidth}\raggedleft
Balanced accuracy (\%)
\end{minipage} & \begin{minipage}[b]{\linewidth}\raggedleft
Mean capped delay
\end{minipage} \\
\midrule\noalign{}
\endhead
\bottomrule\noalign{}
\endlastfoot
Independent & TAN & 0 & 1 & 99.9512 & 0.560547 \\
Independent & Adaptive & 0 & 137 & 93.3105 & 0.103516 \\
AR(0.9) & TAN & 62 & 6 & 96.6797 & 0.930664 \\
AR(0.9) & Adaptive & 7 & 40 & 97.7051 & 0.301758 \\
\end{longtable}

The primary TAN-minus-adaptive AR balanced-accuracy difference is
-1.0254 percentage points, with a predeclared 97.5\% paired percentile
interval {[}-2.0020, -0.0488{]}. The change in advantage from
independent to AR noise is -7.6660 points, interval {[}-9.0820,
-6.2988{]}. Both go against the prespecified superiority direction.

TAN still has fewer catch false alarms under AR noise, but misses more
steps and responds more slowly. This is an operating tradeoff, not
absence of useful behavior. The archived checker reports verification of
10,244 trajectories and all spike decisions; it is a separately
implemented computational check, not external replication. TAN stores
ten input coordinates plus membrane, whereas the adaptive comparator
stores two scalars. The experiment therefore does not establish
resource-normalized superiority for TAN.

\subsection{Modified carriers and exploratory
constructions}\label{modified-carriers-and-exploratory-constructions}

The native eight-stage series illustrates how a conditional rescue can
require a changed contract. A longer gate time constant reduces
sensor-only composition error from 0.270803 to 0.008093 on 4,096
held-out histories. However, the successful system also uses
value-dependent gate drive, a nonzero query content probe, two analog
samples, a known-equation nonlinear inverse, fixed known keys/timing,
and targets restricted to \([-0.5,0.5]^2\). It is neither original P1B
success nor general variable-key retrieval {[}E5{]}.

The interpolation control makes another useful separation. Two real
payload registers can encode the line through two supplied key--value
pairs, enabling ideal querying without retaining separate per-key
records after the second write. Its recorded clean composition error is
approximately \(5.08\times10^{-16}\) on 512 histories. This is classical
interpolation with an explicit controller, separate ports, division, and
atomic overwrite, not a new associative-memory identity or a
demonstrated native neuron. Finite precision, leakage, and resource
costs materially alter the result {[}E6{]}.

Exploratory Chat-sandbox packages address related but different
constructions. The bounded write/hold/query circuit adds state and
control resources; a same-resource EMA performs at least as well in its
stated small comparison. Moment-based latest-value schemes rely on
address promises, conditioning, admissibility tests, or refusal. They
can preserve a response at a noisy observed address while returning an
incorrect value at the true key. These are informative conceptual
controls, not replications of original TAN, and have not been
independently rerun for this manuscript {[}E7{]}.

\section{Discussion}\label{discussion}

The evidence supports an operational distinction between access to
information and the mere presence of information in state. Native
attention can be sensitive to higher moments of the input distribution
while remaining permutation-invariant at a fixed query. A common reset
can remove order-dependent membrane information before that pooled query
occurs. Conversely, static metric addressing succeeds when the candidate
array and bound values are granted. There is no contradiction: the tasks
expose different resources and ask different questions.

The training control adds a second distinction, between a map's
information content and the performance of a particular fitted decoder.
Improving the training distribution can improve a decoder on a fixed
representation. Such improvement does not invert an exact collision and
does not certify sufficient representation. Likewise, the existence of
one collision family does not explain all empirical error. Reporting
both findings prevents either a purely architectural or a purely
optimization-based explanation from outrunning the evidence.

The adaptive comparison shows why a named baseline and a declared
matching rule are necessary. A large independent-noise advantage can
reverse when temporal correlation changes, even without retuning. Fewer
false alarms alone do not establish better detection, just as higher
late activity alone does not establish better information transfer.
Current integrals and logical state counts are useful accounting
quantities, but they are not measurements of energy, biological cost, or
hardware throughput.

The project-specific contribution of this synthesis is the integration
of these verified boundaries and interventions. Exponential pooling,
nearest-neighbor addressing, permutation invariance, interpolation, and
adaptation are established ideas. The evidence does not justify a claim
of a breakthrough, universal task superiority, a minimal neural state
theorem, or a new general theory of attention.

\section{Limitations and
reproducibility}\label{limitations-and-reproducibility}

This paper is retrospective. Protocols described as frozen were saved
before their respective recorded experiments; that does not make the
present cross-study synthesis prospectively registered. Different TAN
variants, normalization rules, tasks, readouts, and resource budgets
cannot be pooled into a common success percentage. The clean-only
control has three fixed historical seeds and conditional test
resampling, not a population-level training study. The collision proof
applies to its supported box and continuation, not arbitrary noise or
all future inputs. Static addressing assumes available stored
candidates. The adaptive comparison covers one calibration and one
correlation change.

The public archive is screened rather than a byte-identical image of
every original machine file. It retains numeric data and records public
text/path transformations. Historical hashes are distinct from current
public-file hashes. Some reproduction guards require the pristine source
or documented adaptation in a disposable copy. Download verification and
source syntax checks are not experimental replication. No new model fit,
simulation sweep, or scientific experiment was run to write this
manuscript.

Readers can inspect the 680-file snapshot, per-claim ledger, archived
protocols and saved outputs. The accompanying validation script checks
source hashes, extracts selected headline numbers directly from saved
CSV/JSON, and validates the original snapshot inventory without
training. The build script produces the PDF and LaTeX from the canonical
Markdown manuscript. Its output should be interpreted as publication
consistency checking. Independent human review, independent execution in
a specified environment, wider task testing, and physical implementation
remain future work.

\section{Conclusion}\label{conclusion}

TAN provides a concrete test bed for separating temporal sensitivity,
static addressing, binding, and recoverable output. The positive
static-address evidence, native response collision, modest
training-composition gain, and reversed correlated-noise comparison
jointly support a resource- and task-specific account. They do not
support treating attention-like weighting as sufficient for reliable
memory. The useful next step is to test a sharply specified mechanism
and output contract, while preserving both the positive results and the
limitations established here.

\section{References}\label{references}

\begin{enumerate}
\def\labelenumi{\arabic{enumi}.}
\tightlist
\item
  Vaswani, A., et al.~(2017). \emph{Attention Is All You Need}. Advances
  in Neural Information Processing Systems 30.
  \href{https://arxiv.org/abs/1706.03762}{Primary paper}.
\item
  Zaheer, M., et al.~(2017). \emph{Deep Sets}. Advances in Neural
  Information Processing Systems 30.
  \href{https://arxiv.org/abs/1703.06114}{Primary paper}.
\item
  Benda, J., Longtin, A., and Maler, L. (2005). Spike-frequency
  adaptation separates transient communication signals from background
  oscillations. \emph{Journal of Neuroscience}, 25(9), 2312--2321. DOI:
  \href{https://doi.org/10.1523/JNEUROSCI.4795-04.2005}{10.1523/JNEUROSCI.4795-04.2005}.
  \href{https://pmc.ncbi.nlm.nih.gov/articles/PMC6726095/}{Primary full
  text}.
\end{enumerate}

\section{Archived evidence
references}\label{archived-evidence-references}

The identifiers below refer to project artifacts rather than externally
published, peer-reviewed papers. Exact paths, hashes, and scope are
provided in \texttt{evidence\_ledger.json} and \texttt{EVIDENCE.md}
beside this manuscript.

\begin{itemize}
\tightlist
\item
  \textbf{E1.} P1A Rev2 condition summary; metric-address proposition;
  v8 P1A evidence review and saved-ledger checks.
\item
  \textbf{E2.} v8 original-project reassessment report, collision
  results, exact theory certificate and scope review.
\item
  \textbf{E3.} Local clean-training frozen protocol, report, results and
  bundled stabilized-source implementation.
\item
  \textbf{E4.} v11 gate-mechanism report, metric summary and frozen
  protocol.
\item
  \textbf{E5.} Native eight-stage experiment-series summary and final
  timescale report.
\item
  \textbf{E6.} Interpolation-control report and subsequent
  finite-precision binding report.
\item
  \textbf{E7.} Exploratory latest-value, range, moment, minimal and
  write/hold/query package documentation.
\end{itemize}

\end{document}
